\section{最小反例}\label{sec:smallest-counterexample}

本节证明定理~\ref{thm:smallest-counterexample}。证明包含逻辑上相互独立的两部分。
首先，用精确矩阵不等式认证一个显式的 32 阶赋号。其次，用完整有限枚举排除
32 以下的每个允许阶数。第一部分足够简短，直接在此给出；第二部分的商集与
证书覆盖性在附录~\ref{app:orbit-completeness} 中证明。

\subsection{32 阶赋号}

令 \(n=32\)，取 Hamilton 和乐 \(\alpha=+1\)，并把三角形通量字

\begin{equation}
\label{eq:order-thirty-two-triangle-word}
\tau=(+,+,-,+,-,-,+,-)\text{，}
\end{equation}

重复四次。在 Hamilton 规范下，所有步长一边均为正，而步长二边
\(\{i,i+2\}\) 的符号为 \(\tau_i\)。所得四边形通量为

\begin{equation}
\label{eq:order-thirty-two-flux-word}
Q=(+,-,-,-)^{8}.
\end{equation}

以 \(A_{32}\) 表示所得符号邻接矩阵。

\begin{proposition}
\label{prop:order-thirty-two-witness}
矩阵 \(A_{32}\) 满足

\begin{equation}
\label{eq:order-thirty-two-certificate}
\rho(A_{32})^{2} < \frac{1561}{200} < \rho_-(32)^{2}.
\end{equation}

\end{proposition}
\begin{proof}
第~\ref{sec:periodic-construction} 节的八点 Floquet 计算给出式~
\eqref{eq:floquet-polynomial} 中的显式多项式 \(P(y,c)\)。令
\(B=\frac{1561}{200}\)、\(u=y-B\)、\(t=2-c\)。式~
\eqref{eq:rational-positive-expansion} 把 \(P(B+u,2-t)\) 展开为系数均非负的
多项式，其常数项为 \(\frac{84332641}{1600000000}>0\)。因此，对每个
\(y\geq B\) 和 \(c\in[-2,2]\)，都有 \(P(y,c)>0\)。单位圆纤维的每个平方
特征值都是 \(P\) 的非负根，故每个纤维谱半径的平方都严格小于 \(B\)。在
\(L=4\)、\(\alpha=+1\) 时，有限分解式~
\eqref{eq:target-finite-decomposition} 给出
\(\rho(A_{32})^2<\frac{1561}{200}\)。

对另一个不等式，式~\eqref{eq:conjectured-threshold} 给出

\begin{equation}
\label{eq:order-thirty-two-threshold}
\begin{aligned}
&\rho_-(32)^{2} \\
&=4+\sqrt{2+\sqrt{2}}+\sqrt{2+\sqrt{2+\sqrt{2}}}
\end{aligned}
\text{。}
\end{equation}

该代数数是下列多项式在区间
\((\frac{7809}{1000},\frac{781}{100})\) 中唯一的实根：

\begin{equation}
\label{eq:order-thirty-two-polynomial}
\begin{aligned}
&X^{8}-32X^{7}+432X^{6}-3216X^{5}+14456X^{4} \\
&\quad -40224X^{3}+67736X^{2}-63184X+25022
\end{aligned}
\text{。}
\end{equation}

Sturm 序列在两个端点的变号数之差为一，且把嵌套根式精确代入即可识别该
孤立根。由于

\begin{equation}
\frac{1561}{200} < \frac{7809}{1000},
\end{equation}

式~\eqref{eq:order-thirty-two-certificate} 中第二个不等式成立。

该证书不依赖于规范。例如，选择步长一边符号

\begin{equation}
(+,-,+,-,-,+,-,+)^{4}\text{，}
\end{equation}

并由 \(b_i=\tau_i a_i a_{i+1}\) 求解，可得步长二边符号

\begin{equation}
(-,-,+,+,+,+,-,-)^{4}.
\end{equation}

对角切换向量 \((+,+,-,-,+,-,-,+)^4\) 把所得矩阵共轭到 \(A_{32}\)。
由此得到一个从通量数据而非所存边列表出发的独立重构。

\end{proof}
\subsection{穷举排除 32 以下阶数}

猜想的定义域恰为偶数 \(n\geq 8\)，故 32 以下的允许阶数是

\begin{equation}
\label{eq:excluded-orders}
8,10,12,14,16,18,20,22,24,26,28,30.
\end{equation}

\begin{proposition}[有限排除]
\label{prop:finite-exclusion}
对于式~\eqref{eq:excluded-orders} 中的每个阶数以及 \(C_n(1,2)\) 的每个赋号
\(\sigma\)，都有

\begin{equation}
\label{eq:finite-exclusion-bound}
\rho(A_\sigma)\geq \rho_-(n).
\end{equation}

这是一条有限计算机辅助命题。

\end{proposition}
\begin{proof}
用圈坐标 \((\tau,\alpha)\) 表示切换类。全局边符号取反和二面体作用只在其
谱不变性得到证明后使用。对 \(n\leq 20\)，直接枚举全部 \(2^{n+1}\) 个类。
对 \(n=22\) 以及正式计算阶数 \(24,26,28,30\)，把合法四边形字划分为二面体
手镯，并保留两种和乐。附录~\ref{app:orbit-completeness} 证明轨道大小之和为
\(2^{n+1}\)，并给出独立的 Burnside 计数。

对特殊最优相位以外的每个代表元，枚举都构造一个非零有理向量 \(v\)，并精确
验证

\begin{equation}
\label{eq:rayleigh-certificate}
(v^T A_\sigma^{2} v)/(v^T v) \geq  \rho_-(n)^{2}.
\end{equation}

式~\eqref{eq:rayleigh-certificate} 中的三角阈值按实代数数比较；浮点特征值只
用于提出候选 \(v\)。特殊类则用其已知精确谱检查。因此每个切换类都满足
式~\eqref{eq:finite-exclusion-bound}。表~\ref{tab:minimality-counts} 的状态数
和证书总数与完整商计数一致，且不存在未判定的代表元。
紧凑归档保存计数、有序判定摘要、检查点链和未认证状态为零的汇总；各个有理
向量由确定性程序重新生成，而不逐项存储。

\begin{table}[tbp]
\centering
\caption{有限最小性枚举与非最优代表元的精确证书。}
\label{tab:minimality-counts}
\begingroup
\renewcommand{\arraystretch}{1.12}
\begin{tabular}{lll}
\toprule
\(n\) & 所代表的切换类 & 非最优代表元的精确证书 \\
\midrule
8 & 512 & 510 \\
10 & 2,048 & 2,046 \\
12 & 8,192 & 8,190 \\
14 & 32,768 & 32,766 \\
16 & 131,072 & 131,070 \\
18 & 524,288 & 524,286 \\
20 & 2,097,152 & 2,097,150 \\
22 & 8,388,608 & 97,467 个商证书 \\
24 & 33,554,432 & 353,811 个商证书 \\
26 & 134,217,728 & 1,299,063 个商证书 \\
28 & 536,870,912 & 4,810,471 个商证书 \\
30 & 2,147,483,648 & 17,929,599 个商证书 \\
\bottomrule
\end{tabular}
\endgroup
\end{table}

最后一列计数的是谱代表元，而不是切换类。附录~
\ref{app:orbit-completeness} 中的商重数恰好解释二者之差。

\end{proof}
\subsection{定理~\ref{thm:smallest-counterexample} 的证明}

命题~\ref{prop:finite-exclusion} 排除了 32 以下的每个允许阶数。命题~
\ref{prop:order-thirty-two-witness} 给出 32 阶的严格反例。由于定义域由不小于
八的偶数组成，30 与 32 之间没有允许阶数。因此 32 是最小反例阶数。

该定理既不声称 32 阶反例唯一，也不声称猜想在 32 以上每个偶数阶都失效。
